% =========================================================
% Proof Stub: Cauchy Property is Inherited by Subsequences
% Source: volume-iii/analysis/sequences/notes/sequences/notes-subsequence-toolkit.tex
% Mode: standard
% =========================================================

\clearpage
\phantomsection
\hypertarget{proof-RA-ABB-P-cauchy-inherited}{}

\subsubsection[Cauchy Property is Inherited by Subsequences]{Proof --- Cauchy Property is Inherited by Subsequences}

\bigskip

\noindent
\textbf{Source.}
See \texttt{notes-subsequence-toolkit.tex}.

\vspace{0.75em}

\noindent
\textbf{Goal.}
If $(a_n)$ is Cauchy, then every subsequence $(a_{n_k})$ is Cauchy.

\vspace{0.75em}

\noindent
\textbf{Logical form.}
\[
  (a_n) \text{ Cauchy} \Rightarrow (a_{n_k}) \text{ Cauchy for every strictly increasing } (n_k).
\]

\vspace{0.75em}

\noindent
\textbf{Proof strategy.}
Given $\varepsilon > 0$, get $N$ from the Cauchy condition on $(a_n)$. For $j, k \ge N_0$ where $n_{N_0} \ge N$, both $n_j, n_k \ge N$, so $|a_{n_j} - a_{n_k}| < \varepsilon$.

\vspace{0.75em}

\noindent
\textbf{Proof.}
\begin{proof}
\Claim If $(a_n)$ is Cauchy, then every subsequence $(a_{n_k})$ is Cauchy.

\Given 

\Goal 

\Strategy Given $\varepsilon > 0$, get $N$ from the Cauchy condition on $(a_n)$. For $j, k \ge N_0$ where $n_{N_0} \ge N$, both $n_j, n_k \ge N$, so $|a_{n_j} - a_{n_k}| < \varepsilon$.

\medskip

% --- given ε > 0, get N: m,n ≥ N ⇒ |a_m - a_n| < ε ---
% --- by Index Growth, ∃ N_0 with n_{N_0} ≥ N ---
% --- for j, k ≥ N_0: n_j, n_k ≥ n_{N_0} ≥ N, so |a_{n_j} - a_{n_k}| < ε ---

\end{proof}

\begin{remark*}[Proof shape]
Index Growth ensures subsequence indices eventually exceed any threshold; past that point the Cauchy condition applies.
\end{remark*}

\begin{remark*}[Dependencies]
\begin{itemize}
  \item Cauchy sequence definition.
  \item Index Growth Lemma: $n_k \to \infty$.
\end{itemize}
\end{remark*}
